Quantum full-state revivals---the exact, periodic reconstruction of an
initial wave packet---are a surprising yet fundamental phenomenon that
occurs in discrete-time quantum walks (DTQW) on symmetric topologies.
These revivals have practical implications for quantum memory \cite{Chandrashekar2015} and 
quantum state routing \cite{Steinfurth2026}, and suggest applications in quantum frequency division, 
since the walker returns to its initial state at integer multiples of the revival period.
Intrinsic (non-contrived) revivals are rare and tend to appear on less
complex topologies with smaller numbers of nodes due to the requirement
that all eigenphases of the system must be exact rational fractions of
$2\pi$. For a full-state revival, every phase step $\theta_j$ of the
unitary step operator $U$ must satisfy
$\theta_j = 2\pi \cdot \frac{m_j}{N}$ for an integer index $m_j$ and
period $N$. As a topology scales in size or complexity, the discrete
momentum spectrum becomes increasingly dense. As the walker progresses
over time, the system accumulates a set of relative phase shifts that
cannot simultaneously evaluate to rational fractions of $2\pi$ for any
common integer step $N$---that is, the corresponding eigenvalues cannot
all be roots of unity. As a result, the wave packet exhibits
irreversible spatial dispersion across the grid, making revivals
impossible.
